%%%PAGE 104 rot=0 legibility=fair summary=牛顿环（薄膜干涉装置与条纹特点）；圆盘衍射（泊松亮斑）；圆孔衍射；单缝衍射；光的干涉；单缝衍射与双缝干涉对比表
%%%BLOCK id=104-01 ch=15 sec=薄膜干涉·牛顿环 kind=knowledge alt=none cont=none y=0.05-0.23
\textbf{牛顿环}

\unclear{多个}透镜接触而干涉

%FIG p104_a 0.045 0.124 0.185 0.172
\notefig[0.28]{figures/p104_a.png}
%FIG p104_b 0.285 0.113 0.435 0.226
\notefig[0.3]{figures/p104_b.png}

$\Delta r=2d$ % CHECK: 原稿字形似"Y"(该书写者的r常写成γ/Y形，参见106页r_2-r_1)，按光程差Δr记

\begin{tabular}{@{}c@{\qquad}c@{}}
单色光 & 白光 \\
$\downarrow$ & $\downarrow$ \\
中暗边亮 & 彩 \\
\end{tabular}
\quad 越近越宽. % 原稿"中暗边亮"与"彩"之间有"/"分隔，单色光、白光的箭头分别指向二者
%%%END
%%%BLOCK id=104-02 ch=15 sec=光的衍射·圆盘衍射 kind=knowledge alt=none cont=none y=0.23-0.37
\textbf{圆盘衍射（泊松亮斑）}\quad 光照到半径很小的圆盘上.

%FIG p104_c 0.213 0.2555 0.378 0.352
\notefig[0.28]{figures/p104_c.png}

\begin{itemize}
\item 中心亮斑
\item 阴影外有\unclear{大}不等间距明暗相间的圆环. % CHECK: "大"字疑为"大小"或指示箭头
\end{itemize}
%%%END
%%%BLOCK id=104-03 ch=15 sec=光的衍射·圆孔衍射 kind=knowledge alt=none cont=none y=0.37-0.57
\textbf{圆孔衍射}

%FIG p104_d 0.225 0.400 0.328 0.466
\notefig[0.25]{figures/p104_d.png}

\textbf{单色光}
\begin{itemize}
\item 中央大亮斑\quad 明暗相间的同心圆环
\item 越往外\quad 条亮纹亮度减弱，宽度减小. % CHECK: 原文"条亮纹"，疑为"亮条纹"
\item 亮环或暗环间的距离随圆孔半径增加而减小,
\end{itemize}

\textbf{白光}\quad 中白边彩\quad 中央大且亮的白色亮斑\ 周围是不等距彩色的同心圆环.
%%%END
%%%BLOCK id=104-04 ch=15 sec=光的衍射·单缝衍射 kind=knowledge alt=none cont=none y=0.57-0.80
\textbf{单缝衍射.}

%FIG p104_e 0.12 0.60 0.292 0.692
\notefig[0.3]{figures/p104_e.png}

\textbf{单色光}
\begin{itemize}
\item 中央为亮且宽的条纹，两侧为明暗相间的条纹
\item 越靠近两侧，亮条纹亮度减弱，宽度越小.
\end{itemize}

\textbf{白光、}
\begin{itemize}
\item 中央为亮且宽的白色条纹，两侧为宽度逐渐变窄的彩色条纹.
\item 最靠近中央的是紫光，离中央最远的是红光.
\end{itemize}
%%%END
%%%BLOCK id=104-05 ch=15 sec=光的干涉·双缝干涉 kind=knowledge alt=none cont=none y=0.80-0.92
\textbf{光的干涉}

%FIG p104_f 0.128 0.835 0.31 0.898
\notefig[0.3]{figures/p104_f.png}

条纹宽度相等、等距、亮度相同.
%%%END
%%%BLOCK id=104-06 ch=15 sec=干涉与衍射比较 kind=summary alt=none cont=none y=0.90-1.00
\begin{center}
\begin{tabular}{ll}
\toprule
单缝衍射 & 双缝干涉 \\
\midrule
条纹宽度不等，中央最宽 & 条纹宽度相等. \\
相邻条纹间距不等 & 相邻条纹等间距 \\
中央条纹最亮，两边变暗 & 条纹清晰，亮度基本相等 \\
\bottomrule
\end{tabular}
\end{center}

本质都是光的叠加 % 原稿此句位于表格左侧；页面底边截断，表格最后一行右栏"亮度基本相等"为两行书写
%%%END
%%%PAGEEND 104
